\section{Discussion}
\subsection{Larger rebound point estimates occur in high, green environments}
The environmental meaning of rebound comes from within-stratum changes rather than the island-wide mean alone. Between Inter-Hurricanes and Post-Fiona, low-elevation low- and high-greenness cells increased by 0.39 and 1.07 percentage points, compared with 6.97 points in the high-elevation, high-greenness cell. The latter also had a high initial reporting level, so its larger absolute increment must be interpreted against that baseline. Both lowland cells retained later-period point estimates below their initial values, suggesting that rebound did not take the form of a uniform return to earlier reporting across environments. Spatially clustered intervals for every Post-Fiona versus initial contrast span zero. The evidence establishes environmental differences in reporting levels more clearly than a categorical separation into recovered and unrecovered environments.

This interpretation also depends on which environments enter the analysis. High-elevation NDVI missingness reached 49.6\%, and high-elevation checklists constituted 11.2\% of missing records but only 3.4\% of complete records. The missing-record reporting odds ratio declined from 1.86 to 1.03 after elevation and period adjustment, identifying an environmental source of aggregate reporting differences. Standardization improves period comparability by holding within-cell covariate distributions constant, while recovering vegetation information at missing-record locations would broaden environmental coverage. Together, these steps help locate rebound without treating changes in sampling weights as changes shared by all sites. This concern is consistent with the interacting sampling biases identified by Backstrom et al.~\cite{backstrom2025}; their results support checking environmental representation alongside survey effort.

Elevation and greenness associations are compatible with the ecological context of forest damage and bird responses following Puerto Rican hurricanes~\cite{uriarte2019,wunderle1995}. Howe et al.~\cite{howe2025} used LiDAR and NDVI to describe complementary structural and greenness responses in mangroves. Their distinction matters here because NDVI alone cannot establish canopy recovery or floral-resource availability. Rivera-Milán et al.~\cite{rivera2025} found prolonged post-hurricane depletion in the Puerto Rico Plain Pigeon, further showing that a common disturbance does not imply a common recovery trajectory across bird taxa. These comparisons provide ecological context rather than independent confirmation of the Emerald pattern. Elevation also represents climate, habitat and accessibility, and the three periods combine two storms and elapsed time without an undisturbed counterfactual, preventing attribution of the period contrast specifically to Fiona.

\subsection{High reporting and persistence coincide in uplands but require distinction within lowlands}
High-elevation, high-greenness environments combine high reporting with high occupancy persistence probability, giving directionally consistent environmental patterns. Correcting observation alignment retained the associations of elevation and greenness with lower local extinction while reducing the absolute initial-period detection coefficient by approximately 90\%. This retains support for linking high reporting in higher, greener settings with occupancy persistence and weakens the earlier strong period-detection contrast. Because the reporting and occupancy models share observations, their agreement compares different response definitions rather than providing independent biological validation.

Lowland results refine this interpretation. The later reporting prediction for the low-elevation, high-greenness cell was 5.17\%, below its initial 6.27\%, yet mean occupancy persistence probability at corresponding sites was 71.3\%, compared with 24.2\% in the low-elevation, low-greenness cell. Incomplete reporting rebound can therefore coexist with relatively strong occupancy persistence probability. The former describes checklist-level reporting trajectories and the latter retention of previously occupied sites. Monitoring must consequently revisit both previously reporting and previously nonreporting sites, estimating state transitions under imperfect detection. Treating all lowlands as sharing a recovery state would obscure their greenness-related differences.

At a fixed elevation of 100~m, changing NDVI from 0.30 to 0.76 corresponded to a 42.2-point persistence contrast but only a 0.7-point colonization contrast with an interval spanning zero. The model supplies clearer support for greenness-associated retention than for greenness-associated entry into occupancy. Under equal visit counts, repeat reports occurred in 1/9 and 15/28 low- and high-greenness lowland pairs, providing a directionally similar observation pattern, although these small-sample histories combine detection and state changes. Greener lowlands therefore merit targeted occupancy surveys, with less-green sites retained as contrasts to distinguish lower occupancy probabilities from inadequate detection.

These probabilities describe between-site associations at the model's scale. Cross-period mean NDVI does not establish the temporal sequence of vegetation recovery within sites, and repeat visits within multi-year primary periods rely on closure. The results identify where to examine persistence; determining whether vegetation change improves persistence requires simultaneous state and vegetation observations within shorter, defensible survey windows. Gordon et al.~\cite{gordon2024} linked colonization and extinction to time since management disturbance using repeated point counts. For the Emerald, recording disturbance and vegetation histories at fixed grids would similarly allow habitat change to be related to separate transition processes.

\subsection{Monitoring should retain upland persistence sites and lowland environmental contrasts}
Only 3.3\% of lowland checklists reported the species, yet lowlands contained 60.7\% of complete-sample positive reports. Lowlands are therefore both difficult reporting environments and an indispensable part of persistence monitoring. Random forest improved AP at low and middle elevations without a comparable high-elevation gain, locating its value in finding sparse reports among many nonreports. Within the global top decile, lowland positive coverage increased from m4's 26.4\% to 34.6\%, while upland coverage declined. Better prediction does not automatically yield more balanced environmental monitoring; ranking must serve the intended environmental contrasts.

Equal-budget comparison quantifies this choice. Global random-forest ranking selected 57.9\% of high-elevation checklists but 6.9\% of low-elevation checklists. Allocating approximately 10\% within each band increased lowland positive coverage to 45.6\% while capturing 388 fewer positives overall. Global ranking favors aggregate hits, whereas comparing occupancy persistence across lowland conditions requires explicit lowland allocation. This comparison concerns historical checklist screening, not the number of independent survey sites, actual revisit costs or improved precision of persistence estimates.

Monitoring can therefore combine three complementary components. First, fixed upland sites can supply detection histories for estimating occupancy transitions in high-reporting environments, with added coverage of sparse conditions such as less-green uplands. Second, lowland coverage should include low, middle and high greenness to compare reporting changes and occupancy persistence across greenness conditions. Third, spatially evaluated ranking can identify supplementary candidates within these strata while fixed sites and nonreporting sites remain in the survey. Candidate screening would then complement complete state-transition observations. Repeat visits should record duration, distance, concurrent vegetation and available floral-resource information to distinguish environmental change from detection opportunity. The fixed routes and interval recordings in the bird-monitoring system of Ovaskainen et al.~\cite{ovaskainen2026} provide a recent example of stabilizing observation conditions alongside predictive modelling. Here, that principle supports maintaining comparable repeat visits even when model rankings change.

\subsection{Spatial and temporal validation delimit the use of revisit clues}
Prediction experiments support these monitoring roles to different degrees across applications. Random forest's AP advantage over m4 narrowed from 0.143 under random splitting to 0.033 under spatial splitting. High scores within familiar location mixtures therefore overstate the corresponding gain in held-out geographic blocks. Increasing block width further reduced some flexible-model advantages. Spatial validation helps assess the domain over which candidate ranking remains informative and supports retaining environmental coverage as surveys expand~\cite{roberts2017,valavi2019,ploton2020}. The current island-internal evaluation without distance buffers does not establish performance in every unsampled environment. Koldasbayeva and Zaytsev~\cite{spacetime2025} likewise emphasize blocked evaluation and out-of-time testing for distribution prediction. Their results motivate matching validation to the intended revisit domain; they do not establish that the block sizes used here remove all spatial dependence.

In the Post-Fiona temporal holdout, random forest increased top-decile recall from m4's 53.28\% to 59.80\%, showing that earlier relationships still helped identify later reports. This supports using existing records to locate follow-up opportunities but does not separately establish temporal gains in each lowland greenness cell. Spatial stratification identifies the environments that benefit, whereas temporal holdout assesses whether aggregate reporting information carries into a later period. Their joint interpretation helps plan revisits without replacing repeated same-site detection records and occupancy-transition analysis. Lovell et al.~\cite{lovell2025} found that spatial and temporal environmental effects on occupancy and abundance differed across parts of a butterfly's range. This offers a methodological parallel for testing temporal performance within environmental strata rather than assuming that an overall holdout gain applies uniformly.

The remaining experiments specify how prediction tools serve the environmental question. EVI2 results show that ranking gains are not confined to NDVI; calibration separates candidate ranking from expected-report estimation; effort responses and importance analyses emphasize comparable survey conditions. Similar spatial performance of m3 and m4 shows that an interaction describing period--elevation patterns need not improve new-site ranking, reinforcing separate evidential requirements for explanation and prediction~\cite{elith2009}. Neural training-source overlap weakens the use of its high scores to justify new monitoring locations. Its curves and ablations describe fitted information dependencies but do not strengthen process conclusions about persistence. Maps and curves should ultimately identify environmental contrasts for revisits, with repeated surveys and detection models estimating whether occupied grids remain occupied in later periods.

\section{Conclusions}
Puerto Rican Emerald reporting rebound is structured by elevation and greenness. High-elevation, high-greenness environments combine larger absolute rebound point estimates with high reporting levels, while recovery relative to initial levels within each stratum requires further repeated observation. These environments also have high modeled persistence. Greener lowlands retain a stronger persistence association, showing that incomplete lowland reporting rebound cannot be equated with nonoccupancy at all previously occupied grids. Prediction gains chiefly identify sparse low- and middle-elevation reports, and spatial and temporal validation delimit their use as revisit clues. Monitoring should combine fixed upland repeat sites, lowland greenness contrasts and within-stratum candidate screening to examine reporting changes and repeat records separately from occupancy persistence and colonization estimated under imperfect detection.

\FloatBarrier
